\documentclass[11pt]{article}
\usepackage[margin=0.9in]{geometry}
\usepackage{amsmath,amssymb,amsthm,booktabs,mathrsfs}
\usepackage[T1]{fontenc}
\usepackage{lmodern,microtype}
\usepackage[colorlinks=true,linkcolor=blue,urlcolor=blue]{hyperref}
\newtheorem{theorem}{Theorem}[section]
\newtheorem{lemma}[theorem]{Lemma}
\newtheorem{proposition}[theorem]{Proposition}
\newtheorem{corollary}[theorem]{Corollary}
\theoremstyle{remark}
\newtheorem{remark}[theorem]{Remark}
\title{Toward a four-dimensional Yang--Mills construction\\
\large Notebook 34: ghost asymptotics and the physical determinant}
\author{Research working document}
\date{17 October 2026}
\begin{document}
\maketitle
\begin{abstract}
We determine the complete based-ghost contribution to the
actual leading determinant's Maxwell coefficient. An exact
returning-walk area calculation gives the unpinned lattice
heat coefficient. Counting the root residues excluded by
the pins controls the difference through a specified
heat-time interval. Through the block diffusion time
\(K^2\), the actual contribution has a fixed finite part
and an \(O(K^{-2})\) error.

The flat pinning gap is of order \(K^{-4}\), rather than
\(K^{-2}\). A pointwise Neumann Green-function estimate
proves the matching lower bound. A fourth area moment
and a killed-return heat bound then control the tail
past \(K^4\). Together they give the entire actual ghost
coefficient \(C_H\log K/(12\pi^2)+O(1)\), uniformly in
the preceding classical depth and the permitted infinite
structural systems.

This is twice the ghost logarithm of the massive reference
used in Notebook 33. The extra logarithm comes from the
interval \(K^2\) to \(K^4\) and must be compensated by the
actual vector determinant. It supplies a quantitative
test of the remaining coupled comparison, not the full
Yang--Mills existence argument.

We also prove the full actual vector short-time bound,
including the composite ordinary constraint. Its flat
derivative has a uniform counting-\(\ell^2\) bound,
and its affine magnetic jets are controlled through
the fixed-ratio hierarchy. A localized trace identity
uses the fixed coarse fiber dimension to avoid a
factor equal to the growing fine block volume.
Together with the ghost estimate, this closes the
actual ultraviolet heat interval.

Finally an exact finite determinant identity separates
the full fine and coarse gauge-orbit density factors
from a residual precision on the physical quotient.
That physical precision has a uniform flat gap of
order \(K^{-2}\). An explicit contour-gauge test shows
why the same gap cannot simply be assigned to the
unprojected contour completion. A paired scalar
reference has an exact infrared cancellation; applying
it to the actual finite identity still requires a
joint infrared limit and uniform curvature-source
estimates for the physical operator.
\end{abstract}

\section*{Inputs, conventions and scope}

The actual construction, fixed-\(K\) local infinite limits
and determinant-density identities are those of Notebooks
32 and 33. The input identities are:
\begin{enumerate}
\item On the normalized output probe
\(w_\mu(n)=t n_\nu H\), \(\mu\ne\nu\), \(\tau(H^2)=1\),
the first two actual stationary jets agree with current
affine links \(U_\mu(x)=\exp(i t x_\nu H/K^2)\).
This holds for every preceding homogeneous classical
action \(S_M\). Only these jets are used.
\item The based generator \(D_b\) has scalar domain zero
at every \(K\)-center, and \(G(t)=D_b(t)^TD_b(t)\).
The Maxwell contribution of its normalized log determinant
is \(2[t^2]\operatorname{Tr}_{0,S}\log G(t)\), where the
trace includes noncenter sites in one centered block
and their adjoint color coordinates.
\item The actual total coefficient splits into this ghost
term, the vector term
\(-[t^2]\operatorname{Tr}_{0,E}\log V_a(t)\), and a
uniformly bounded coarse term. Here
\(V_a=s_M''+D_bD_b^T+aC_K^TC_K\) on the required jets;
the constraint multiplier curvature vanishes through
the same two orders. The fixed \(a\) is the large
positive completion from Notebook 32.
\end{enumerate}

The charge factor is \(C_H=\sum_q q^2\), where \(q\)
runs over the eigenvalues of \(\operatorname{ad}_H\)
on the complexified adjoint space, with multiplicities.
There is no doubling of the real adjoint trace.
For \(SU(N)\) with \(\tau=\operatorname{Tr}_{\rm fund}/N\),
\(C_H=2N^2\). We write \(\alpha_g=C_H/(48\pi^2)\).
All estimates below are in four dimensions.

The proofs have been checked internally against these
inputs. This notebook does not assert a full
endpoint-uniform vector comparison, positive total
drift, a continuum measure, or reconstruction.

